\scanpage{51}
\mspage{8}
% archive sheet number: 50
Il semble plus prudent, dans la définition de
% slanted
\margin{notion : \uline{très rigide}}la rigidité \inserted{sur un \uline{corps}}, d'exiger que \unsure{ses|les|tous les} sous-schém. fermés \inserted{\unsure{relt.|rel.!|rel.}} connexes
de $G$ ait la $=$ propriété de densité. Cela
signifie alors que $G$ a une suite de
composition (sur $k$ \uline{alg. clos}) \deleted{\unsure{tq}} dont les facteurs,
sont $\unsure{Z_{\ell}|Z\ell|Z/\ell}$, $G_m$, $\mathrm{Ab}$, --- et \deleted{cette} \deleted{condition} \inserted{notion}
\inserted{de \uline{très rigide}}, comme on voit facilement, se comporte
bien par limites inductives, est stable
par passage : \unsure{aux} \unsure{ss-schémas|ks-schémas} \unsure{fermés|fermé} simples
sur $k$.

\deleted{\uline{Proposition} $G$ rigide sur $S$ loc. noeth., $u:G\longrightarrow H$ un morphisme de $S$-groupes, où $H$ est un schém. en groupes \inserted{\unsure{comm.}} de type fini sur $S$. Soit $N=\operatorname{Ker}u$. Alors $N$ est universellement ouvert sur $S$.}

\deleted{On peut supposer $S=\operatorname{Spec} V$, $V$ anneau de valuation discrète.}

\uline{Remarque} Même si $A$ est un schém.
\margin{\rule[-3\baselineskip]{0.35pt}{3\baselineskip}\kern2pt\rule[-3\baselineskip]{0.35pt}{3\baselineskip}}abélien sur un anneau de valuation discrète $V$,
$B_1$ et $B_2$ deux sous-schémas \inserted{en groupes finis} de $A$ \deleted{\unsure{plats}} \inserted{\uline{étales}}
sur $V$, de degré $p$, ayant $=$ fibre
en \unsure{le} pt fermé $s$, il est possible qu'on ait $B_1\ne B_2$
